% !TeX root = ../thesis.tex
% ===================== ГЛАВА 3 =====================
\chapter{Задача переноса: ударная волна перед поршнем}

\section{Постановка задачи}

Рассматривается задача о поршне в полупространстве $x\le 0$.
В системе отсчёта, связанной с поршнем, стенка покоится при $x=0$, а слева ($x=-L$) втекает сверхзвуковой поток газа со
скоростью $u_0$.
Скорость потока превышает скорость звука $c=\sqrt{\gamma T}\approx 1{,}291$ (одноатомный газ, $\gamma=5/3$); принято
$u_0=2c\approx 2{,}58$, что отвечает числу Маха набегающего потока $M=u_0/c=2$.
Перед поршнем формируется ударная волна, распространяющаяся в отрицательном направлении.
Эволюция описывается одномерным по пространству и трёхмерным по скоростям уравнением Больцмана
% FIX: в описании задачи в task2/docs/task.txt на строке 14 указано, что задача одномерная
\begin{equation}
  \frac{\partial f}{\partial t}+\xi_x\frac{\partial f}{\partial x}=I(f,f),\qquad f=f(x,\xv,t).
\end{equation}
На стенке задаётся зеркальное (адиабатическое) отражение, слева~--- набегающее максвелловское распределение с дрейфом
$u_0$.

Выбор системы отсчёта, связанной с поршнем, удобен тем, что стенка неподвижна, а расчётная область $[-L,0]$ фиксирована
во времени.
Размер области $L$~--- существенный параметр задачи: фронт ударной волны движется влево со скоростью $V$, и за время
счёта он не должен достигать левой границы; кроме того, область должна вмещать кинетический предвестник~--- сравнительно
быстрые молекулы, уходящие вверх по потоку перед фронтом.
Дискретизация по пространству равномерная, $x_i=-L+(i-\tfrac12)h$ и $h=L/N_x$; по скоростям используется та же сетка с
обрезанием, что и в задачах релаксации. % FIX: так, а в каких пределах меняется i?

\section{Схема расщепления и оператор переноса}

Применяется расщепление Странга по физическим процессам (2-й порядок по времени)
\begin{equation}
  \text{Перенос}\Big(\tfrac{\Delta t}{2}\Big)\;\to\;
  \text{Столкновения}(\Delta t)\;\to\;
  \text{Перенос}\Big(\tfrac{\Delta t}{2}\Big).
\end{equation}
Оператор столкновений совпадает с описанным в главе~1, но в неоднородной задаче выделена лишь поперечная симметрия
(по $\xi_y,\xi_z$): ось $x$ выделена направлением течения, поэтому индекс $\xi_x$ не отражается, что даёт 4-кратную
(а не 8-кратную) связь узлов.
Геометрия столкновений (узлы интерполяции, веса $r_\nu$, модули $|\gv_\nu|$) вычисляется один раз и затем векторно
применяется ко всем пространственным ячейкам.

Оператор переноса $\partial f/\partial t+\xi_x\,\partial f/\partial x=0$ решается монотонной разностной схемой TVD
второго порядка в потоковой форме с ограничителем наклонов MC (monotonic central),
\begin{equation}
  \varphi(\theta)=\max\!\big[0,\ \min(2\theta,\ \tfrac{1+\theta}{2},\ 2)\big],
\end{equation}
для каждого скоростного узла со своим числом Куранта $\nu_a=|\xi_{x}|\,\Delta t/h\le 1$.
% FIX: является ли такое обозначение стандартным? как будто в Пособиях было другое
Граничные условия реализованы через слои фиктивных ячеек: слева~--- набегающий максвелл $f_\infty$ с дрейфом $u_0$,
справа~--- зеркальное отражение (узел $\xi_x\to-\xi_x$), дающее точный отражённый поток.

\section{Теория ударной волны}

Параметры установившейся ударной волны даёт теория Ренкина--Гюгонио.
Для принятых условий число Маха \emph{самого скачка} $M_s\neq M$ набегающего потока и составляет $M_s=3{,}0$; при этом
гидродинамические скачки плотности и температуры равны
% FIX: напомни мне, почему $M_s\neq M$, откуда такие скачки
\begin{equation}
  \frac{n_1}{n_0}=3{,}0,\qquad \frac{T_1}{T_0}\approx 3{,}66,
\end{equation}
а скорость фронта $V=u_0-M_s\,c\approx-1{,}29$ (фронт движется влево). % FIX: почему тут вычитание?
Число Маха скачка связано со скоростью фронта соотношением $M_s=(u_0-V)/c$; самосогласованное решение задачи
поршень--скачок даёт $V=-c$, откуда $M_s=3$. % FIX: то есть ты сначал определил V через M_s, а потом M_s через v?
Эти величины служат ориентиром для проверки результатов кинетического расчёта.

\section{Результаты}

Окончательный расчёт выполнен в конфигурации, обеспечивающей полное развитие ударной волны: расчётная область $L=80$,
число пространственных ячеек $N_x=400$, сетка скоростей с $N=30$ и $\xi_{\mathrm{cut}}=7{,}5$, простое число сетки
Коробова $p=50021$, счёт до $t=40$ ($2000$ шагов).
% FIX: в чём у нас измеряется время? и что значат эти шаги, это расщепление странга?
Размер области выбран из условия $L\gtrsim 75$ (см.\ ниже): при меньших $L$ кинетический предвестник достигает левой
границы и компрессия «застывает» на $n\approx 2{,}3$, не выходя на гидродинамический предел.

На рис.~\ref{fig:t2prof} приведены профили концентрации $n$, скорости $u_x$, температуры $T$ и продольной компоненты
$T_{xx}$ вдоль оси $x$ в несколько моментов времени.
Чётко виден установившийся фронт ударной волны, движущийся влево; за фронтом газ сжат
(концентрация выходит на $n\approx3$) и нагрет ($T\approx3{,}6$), а скорость в системе поршня падает почти до нуля.
% FIX: за фронтом, то есть в сторону 0 или -inf?
На самом фронте заметен немонотонный «перескок» $T_{xx}$~--- характерный кинетический (неравновесный) признак,
отсутствующий в гидродинамическом описании через уравнение Навье-Стокса.
Рис.~\ref{fig:t2front} показывает траекторию фронта $x_c(t)$: после короткого переходного участка она выходит на прямую,
наклон которой даёт скорость фронта $V\approx-1{,}22$.

Эволюция одномерной функции распределения $f_x(\xi_x)$ в середине области (рис.~\ref{fig:t2fx}) иллюстрирует кинетику
ударного перехода: до прихода фронта здесь наблюдается узкий пик набегающего холодного сверхзвукового потока при
$\xi_x\approx u_0$, а после прохождения фронта распределение уширяется и смещается к $\xi_x\approx0$~--- газ
термализуется в «послескачковое» распределение Максвелла.
Тем самым в одном расчёте видны обе стороны ударной волны: сверхзвуковой набегающий поток и дозвуковой сжатый газ за
фронтом.

\begin{figure}[h]
  \centering
  \includegraphics[width=0.95\textwidth]{t2_profiles}
\caption{Профили макропараметров $n,\,u_x,\,T,\,T_{xx}$ вдоль оси $x$ в различные моменты времени (установившаяся
ударная волна, конфигурация $L=80$, $N_x=400$).}
  \label{fig:t2prof}
\end{figure}

\begin{figure}[h]
  \centering
  \includegraphics[width=0.8\textwidth]{t2_front}
\caption{Траектория фронта $x_c(t)$ (линейный участок даёт скорость фронта $V$) и эволюция максимальной концентрации
$n_{\max}(t)$.}
  \label{fig:t2front}
\end{figure}

\begin{figure}[h]
  \centering
  \includegraphics[width=0.8\textwidth]{t2_fx}
\caption{Эволюция $f_x(\xi_x)$ в середине области: переход от холодного набегающего потока (пик при $\xi_x\approx u_0$)
к послескачковому распределению ($\xi_x\approx0$).}
  \label{fig:t2fx}
\end{figure}

\section{Проверка теории Ренкина--Гюгонио}

Сравнение измеренных параметров скачка с точной теорией приведено в табл.~\ref{tab:hugoniot}.
Достигнуто хорошее согласие: скачок плотности $n_1/n_0=2{,}95$ (теория $3{,}0$), скачок температуры $T_1/T_0=3{,}62$
(теория $3{,}66$), скорость фронта $V=-1{,}22$ (теория $-1{,}29$); относительная невязка первого соотношения Гюгонио
$n_0(u_0-V)=n_1(u_1-V)$ составляет $2{,}25\,\%$. % FIX: напомни мне, каков физический смысл этого соотношения?
Масса, импульс и энергия сохраняются с машинной точностью (по построению схемы), а доля пропущенных по условию
положительности точек равна $0{,}53\,\%$~--- в пределах допустимого порога.

\begin{table}[h]
  \centering
\caption{Параметры ударной волны: кинетический расчёт против теории Ренкина--Гюгонио (относительное отклонение вычислено
от неокруглённых значений).}
  \label{tab:hugoniot}
  \begin{tabular}{lccc}
    \hline
    Величина & Расчёт & Теория & Отн.~отклонение\\
    \hline
    Скачок плотности $n_1/n_0$   & 2,954      & 2,999      & 1,5\,\%\\
    Скачок температуры $T_1/T_0$ & 3,617      & 3,663      & 1,3\,\%\\
    Скорость фронта $V$          & $-1{,}221$ & $-1{,}291$ & 5,4\,\%\\
    Число Маха скачка $M_s$      & 2,998      & 3,000      & ---\\
    \hline
  \end{tabular}
\end{table}

Отметим несколько физических особенностей.
% FIX: это строка спряталась между таблицей и графиком и труднозаметна. можно ли её как-то элегантно подвинуть?
Во-первых, выход на гидродинамический скачок $n\to3$ происходит не мгновенно: вблизи стенки сначала формируется
свободномолекулярное удвоение концентрации $n\approx2$ (зеркальное отражение набегающего потока), а затем медленная
компрессия через обратную связь переноса и столкновений доводит концентрацию до $n\approx3$; установившийся режим
достигается к $t\approx10$--$15$.

Во-вторых, ключевую роль играет размер области: за фронтом горячий газ испускает быстрые молекулы вверх по потоку
(кинетический предвестник, вплоть до скоростей $\xi_{\mathrm{cut}}$); если левая граница расположена близко, предвестник
покидает область и компрессия не достигает теоретического значения.
Условие $L\gtrsim\xi_{\mathrm{cut}}\cdot t_{\text{форм}}\approx75$ обеспечивает корректное развитие волны.
В-третьих, несколько большее расхождение по скорости фронта ($5{,}4\,\%$ против $\sim1{,}5\,\%$ по $n$ и $T$) объяснимо:
$V$ извлекается из наклона траектории $x_c(t)$, а фронт имеет конечную кинетическую ширину и чувствителен к порогу его
детектирования. % FIX: что такое кинетическая ширина? чему она равна?
К $t=40$ плато плотности не дошло до $3{,}0$ примерно на $1{,}5\,\%$, что слегка занижает $|V|$.
% FIX: на графике видно, что с ростом времени невязка v растёт. почему так и плохой ли это признак?
Локальные скачки $n$ и $T$ сходятся точнее, тогда как $V$ и интегральная невязка~--- наиболее чувствительные метрики;
% FIX: что здесь является интегральной невязкой?
для кинетического расчёта без подгоночных параметров отклонения в $2$--$5\,\%$~--- хороший результат.
Таким образом, реализованный солвер не только сохраняет инварианты и $H$-функцию, но и количественно воспроизводит
ударно-волновое течение в согласии с теорией.

